\vspace{-2mm}
\section{Related Work}
\label{sec:relatedwork}
\vspace{-1mm}

Our work is related to provenance in relational databases, scientific workflows, and LLM inference. 

\topic{Database Provenance} We compared against database provenance in Section~\ref{subsec:database-prov-relationship}. Here, we provide a more general background. 
Database provenance~\cite{cheney2009provenance,buneman2007provenance,glavic2021data,karvounarakis2010querying} focuses on tracking data lineage to ensure reproducibility. 
Given a SQL query and inputs, provenance investigates: 
(1) the input tuples that produce the output (why-provenance)~\cite{cui2000tracing,buneman2001and,buneman2002propagation}, (2) the specific locations (e.g., table, row, or column) in the input from which elements in the output originate (where-provenance)~\cite{buneman2001and,bhagwat2005annotation}, and (3) the transformations applied to the input to derive the output (how-provenance)~\cite{green2007update,green2007provenance,green2007orchestra,hernandez2021computing}. 
Our setting poses additional challenges:
namely the ambiguity of queries and data,
the black-box nature of LLMs,
and dealing with LLM costs in inferring provenance.
% Compared to database provenance, provenance for question answering with LLMs poses unique challenges: first, natural language questions and text data are far more ambiguous than the deterministic nature of table data and SQL queries. Second, LLMs introduce additional complexity, such as monetary cost, requiring a balance between cost and provenance size. 
% \sys introduces the concept of verifiable provenance for question answering over textual data using LLMs, achieving high accuracy in inferring the true provenance while minimizing cost and ensuring minimality. 




\topic{Provenance in Data Science Pipelines} 
Provenance has been widely studied in scientific workflow systems, which use a dataflow model.
%where the execution order of modules is driven by the flow of data.
VisTrails~\cite{scheidegger2008querying} enhances provenance query efficiency by enabling users to reuse common sub-workflows. %, reducing the need for provenance re-computation. 
Subzero~\cite{wu2013subzero} leverages data locality, where multiple operators in a workflow often share similar input cells. %, and introduces a fine-grained lineage system to improve query performance. 
DfAnalyzer~\cite{silva2018dfanalyzer} tracks provenance at runtime without waiting for workflow completion.%, minimizing query overhead through strategies like parallel data loading and metadata column storage.  
Adriane et al.~\cite{chapman2020capturing} focuses on machine learning workflows, computing provenance to debug model performance. 
%by identifying input training data responsible for output degradation. 
\sys instead focuses on computing verifiable provenance within a single LLM invocation, and extending \sys to workflows would be interesting future work. 

%A promising future direction is to adapt \sys for workflows involving LLM operators, given the rising popularity of agentic LLM-based systems. 
 

%rupprecht2020improving
%chapman2022dpds 
%davidson2008provenance 

\topic{Provenance in LLM Inference}
Overall, LLM provenance has only been studied in limited settings. Prior work, including
%While LLMs have emerged as a promising tool for text analytics, the provenance of LLMs' response has not been extensively studied. 
LLM Attributor~\cite{lee2025llm},  ContextCite~\cite{cohen2024contextcite}, and Hithesh et al~\cite{sankararaman2024provenance} all identify top-k relevant tokens as provenance, using embedding similarities, next-token probabilities, or relevance scores predicted by pretrained models. However, such heuristic methods lack verifiability. % to ensure the provenance that can reproduce the same (or equivalent) answers.   
Notably, these approaches can be integrated into \sys as candidate ranking models in our pruning phase. \trs{Note that LLM Attributor~\cite{lee2025llm} requires access to LLM training data, which is often unavailable.} 
%for closed-source models (e.g., OpenAI) and many open-source models (e.g., LLama3), limiting its practicality.  